\documentclass[11pt,a4paper]{article}

% --- THIẾT LẬP LỀ TRANG & KHỔ GIẤY ---
\usepackage[top=1.8cm,bottom=1.8cm,left=1.8cm,right=1.8cm]{geometry}

% --- HỖ TRỢ ĐA DẠNG TRÌNH BIÊN DỊCH (XeLaTeX / pdfLaTeX / LuaLaTeX) ---
\usepackage{iftex}
\ifXeTeX
  \usepackage{fontspec}
  \usepackage[vietnamese]{babel}
  \defaultfontfeatures{Ligatures=TeX}
\else
  \usepackage[utf8]{inputenc}
  \usepackage[T5,T1]{fontenc}
  \usepackage[vietnamese]{babel}
  \usepackage{mathptmx}
\fi

% --- CÁC GÓI TOÁN HỌC CHUẨN AMS ---
\usepackage{amsmath}
\usepackage{amssymb}
\usepackage{amsfonts}
\usepackage{mathtools}
\usepackage{bm}

% --- BẢNG BIỂU, DANH SÁCH & TRANG TRÍ ---
\usepackage{tabularx}
\usepackage{booktabs}
\usepackage{array}
\usepackage{enumitem}
\usepackage{multicol}
\usepackage{fancyhdr}
\usepackage{xcolor}

% --- THIẾT LẬP HEADER & FOOTER ---
\pagestyle{fancy}
\fancyhf{}
\renewcommand{\headrulewidth}{0pt}
\lfoot{\textit{\small Bộ đề kiểm tra Toán THCS KNTT - Mã: c1\_6\_bai2\_test1}}
\rfoot{\textit{\small Trang \thepage}}

\begin{document}

% ==============================================================================
% PHẦN ĐẦU ĐỀ THI (HEADER TIÊU CHUẨN BỘ GD&ĐT)
% ==============================================================================
\noindent
\begin{tabular*}{\textwidth}{@{\extracolsep{\fill}}lr@{}}
\textbf{SỞ GD\&ĐT ........................} & \textbf{ĐỀ KIỂM TRA ĐỊNH KỲ NĂM HỌC 2025 - 2026} \\
\textbf{TRƯỜNG THCS ........................} & \textbf{Môn: TOÁN - LỚP 6} \\
\textit{(Đề thi chuẩn ma trận 2025)} & \textit{Thời gian làm bài: 45 phút (không kể giao đề)}
\end{tabular*}

\vspace{6pt}
\begin{center}
{\Large\bfseries Toán 6 (KNTT) - Bài 2: Cách ghi số tự nhiên - Bài kiểm tra số 1}\\[4pt]
{\bfseries [ MỨC ĐỘ 1: CƠ BẢN - Mã đề: \texttt{c1\_6\_bai2\_test1} ]}
\end{center}

\vspace{4pt}
\noindent
\rule{\textwidth}{0.8pt}\\[3pt]
Họ và tên thí sinh: \dotfill\ Lớp: \dots\dots\ SBD: \dots\dots\\[-6pt]
\rule{\textwidth}{0.8pt}

\vspace{10pt}

% ==============================================================================
% PHẦN I: TRẮC NGHIỆM NHIỀU PHƯƠNG ÁN (3,0 ĐIỂM)
% ==============================================================================
\noindent
\textbf{\large PHẦN I. Câu trắc nghiệm nhiều phương án lựa chọn (3,0 điểm)}\\
\textit{\small Thí sinh trả lời từ câu 1 đến câu 12. Mỗi câu hỏi thí sinh chỉ chọn một phương án đúng nhất.}
\vspace{6pt}
\noindent\textbf{Câu 1.} [Nhận biết] Trong hệ thập phân, người ta dùng bao nhiêu chữ số để ghi tất cả các số tự nhiên?\par\vspace{2pt}
\noindent\begin{tabular*}{\textwidth}{@{\extracolsep{\fill}}llll@{}}
\textbf{A.} $9$ & \textbf{B.} $10$ & \textbf{C.} $8$ & \textbf{D.} Vô số
\end{tabular*}\par\vspace{5pt}

\noindent\textbf{Câu 2.} [Nhận biết] Chữ số La Mã$\mathrm{X}$biểu diễn số tự nhiên nào trong hệ thập phân?\par\vspace{2pt}
\noindent\begin{tabular*}{\textwidth}{@{\extracolsep{\fill}}llll@{}}
\textbf{A.} $1$ & \textbf{B.} $5$ & \textbf{C.} $10$ & \textbf{D.} $50$
\end{tabular*}\par\vspace{5pt}

\noindent\textbf{Câu 3.} [Nhận biết] Chữ số La Mã$\mathrm{V}$có giá trị tương ứng là:\par\vspace{2pt}
\noindent\begin{tabular*}{\textwidth}{@{\extracolsep{\fill}}llll@{}}
\textbf{A.} $5$ & \textbf{B.} $4$ & \textbf{C.} $6$ & \textbf{D.} $10$
\end{tabular*}\par\vspace{5pt}

\noindent\textbf{Câu 4.} [Nhận biết] Trong số tự nhiên$3\,456$, chữ số$4$nằm ở hàng nào?\par\vspace{2pt}
\noindent\begin{tabular*}{\textwidth}{@{\extracolsep{\fill}}llll@{}}
\textbf{A.} Hàng đơn vị & \textbf{B.} Hàng chục & \textbf{C.} Hàng trăm & \textbf{D.} Hàng nghìn
\end{tabular*}\par\vspace{5pt}

\noindent\textbf{Câu 5.} [Nhận biết] Chữ số La Mã$\mathrm{IV}$biểu thị số nào sau đây?\par\vspace{2pt}
\noindent\begin{tabular*}{\textwidth}{@{\extracolsep{\fill}}llll@{}}
\textbf{A.} $4$ & \textbf{B.} $6$ & \textbf{C.} $14$ & \textbf{D.} $16$
\end{tabular*}\par\vspace{5pt}

\noindent\textbf{Câu 6.} [Nhận biết] Trong số tự nhiên$285$, chữ số hàng chục là:\par\vspace{2pt}
\noindent\begin{tabular*}{\textwidth}{@{\extracolsep{\fill}}llll@{}}
\textbf{A.} $2$ & \textbf{B.} $8$ & \textbf{C.} $5$ & \textbf{D.} $80$
\end{tabular*}\par\vspace{5pt}

\noindent\textbf{Câu 7.} [Nhận biết] Số tự nhiên có hai chữ số$\overline{ab}$($a \ne 0$) được viết thành tổng giá trị các chữ số là:\par\vspace{2pt}
\noindent\begin{tabular*}{\textwidth}{@{\extracolsep{\fill}}llll@{}}
\textbf{A.} $a + b$ & \textbf{B.} $10a + b$ & \textbf{C.} $a + 10b$ & \textbf{D.} $100a + b$
\end{tabular*}\par\vspace{5pt}

\noindent\textbf{Câu 8.} [Nhận biết] Số$9$được biểu diễn bằng chữ số La Mã là:\par\vspace{2pt}
\noindent\begin{tabular*}{\textwidth}{@{\extracolsep{\fill}}llll@{}}
\textbf{A.} $\mathrm{VIIII}$ & \textbf{B.} $\mathrm{XI}$ & \textbf{C.} $\mathrm{IX}$ & \textbf{D.} $\mathrm{VIV}$
\end{tabular*}\par\vspace{5pt}

\noindent\textbf{Câu 9.} [Thông hiểu] Giá trị của chữ số$5$trong số$25\,341$là:\par\vspace{2pt}
\noindent\begin{tabular*}{\textwidth}{@{\extracolsep{\fill}}llll@{}}
\textbf{A.} $5$ & \textbf{B.} $50$ & \textbf{C.} $500$ & \textbf{D.} $5\,000$
\end{tabular*}\par\vspace{5pt}

\noindent\textbf{Câu 10.} [Thông hiểu] Số La Mã$\mathrm{XXIV}$tương ứng với số tự nhiên nào trong hệ thập phân?\par\vspace{2pt}
\noindent\begin{tabular*}{\textwidth}{@{\extracolsep{\fill}}llll@{}}
\textbf{A.} $24$ & \textbf{B.} $26$ & \textbf{C.} $19$ & \textbf{D.} $29$
\end{tabular*}\par\vspace{5pt}

\noindent\textbf{Câu 11.} [Thông hiểu] Số$18$được viết bằng chữ số La Mã là:\par\vspace{2pt}
\noindent\begin{tabular*}{\textwidth}{@{\extracolsep{\fill}}llll@{}}
\textbf{A.} $\mathrm{XIIX}$ & \textbf{B.} $\mathrm{XVIII}$ & \textbf{C.} $\mathrm{XV}$ & \textbf{D.} $\mathrm{XVI}$
\end{tabular*}\par\vspace{5pt}

\noindent\textbf{Câu 12.} [Thông hiểu] Số tự nhiên gồm$7$trăm,$0$chục và$4$đơn vị được viết là:\par\vspace{2pt}
\noindent\begin{tabular*}{\textwidth}{@{\extracolsep{\fill}}llll@{}}
\textbf{A.} $740$ & \textbf{B.} $74$ & \textbf{C.} $704$ & \textbf{D.} $7\,004$
\end{tabular*}\par\vspace{5pt}


% ==============================================================================
% PHẦN II: TRẮC NGHIỆM ĐÚNG / SAI (2,0 ĐIỂM)
% ==============================================================================
\vspace{6pt}
\noindent
\textbf{\large PHẦN II. Câu trắc nghiệm đúng sai (2,0 điểm)}\\
\textit{\small Thí sinh trả lời từ câu 1 đến câu 2. Trong mỗi ý a), b), c), d) ở mỗi câu, thí sinh chọn đúng hoặc sai.}\\
\textit{\small (Quy tắc tính điểm: Đúng 1 ý được 0,1đ; đúng 2 ý được 0,25đ; đúng 3 ý được 0,5đ; đúng 4 ý được 1,0đ).}
\vspace{6pt}
\noindent\textbf{Câu 1.} Cho số tự nhiên$N = 3\,528$. Xét tính đúng hoặc sai của các khẳng định sau:\par\vspace{2pt}
\begin{enumerate}[label=\textbf{\alph*)}, leftmargin=1.2cm, itemsep=2pt, topsep=2pt]
  \item \textnormal{[Thông hiểu]} Chữ số hàng nghìn của số$N$là$3$.
  \item \textnormal{[Thông hiểu]} Giá trị của chữ số$5$trong số$N$là$50$.
  \item \textnormal{[Thông hiểu]} Số$N$được viết thành tổng giá trị các chữ số là:$3 \times 1\,000 + 5 \times 100 + 2 \times 10 + 8$.
  \item \textnormal{[Thông hiểu]} Số các chữ số của số$N$là$5$.
\end{enumerate}
\vspace{6pt}

\noindent\textbf{Câu 2.} Xét các số La Mã biểu diễn các số tự nhiên:\par\vspace{2pt}
\begin{enumerate}[label=\textbf{\alph*)}, leftmargin=1.2cm, itemsep=2pt, topsep=2pt]
  \item \textnormal{[Thông hiểu]} Kí hiệu$\mathrm{VI}$biểu diễn số$6$.
  \item \textnormal{[Thông hiểu]} Kí hiệu$\mathrm{XI}$có giá trị nhỏ hơn kí hiệu$\mathrm{IX}$.
  \item \textnormal{[Thông hiểu]} Số$19$được viết dưới dạng số La Mã là$\mathrm{XIX}$.
  \item \textnormal{[Thông hiểu]} Số La Mã$\mathrm{XXVII}$có giá trị bằng$25$.
\end{enumerate}
\vspace{6pt}


% ==============================================================================
% PHẦN III: TRẮC NGHIỆM TRẢ LỜI NGẮN (2,0 ĐIỂM)
% ==============================================================================
\vspace{6pt}
\noindent
\textbf{\large PHẦN III. Câu trắc nghiệm trả lời ngắn (2,0 điểm)}\\
\textit{\small Thí sinh trả lời từ câu 1 đến câu 4. Mỗi câu trả lời đúng thí sinh được 0,5 điểm.}
\vspace{6pt}
\noindent\textbf{Câu 1.} [Thông hiểu] Tính giá trị của chữ số$7$trong số tự nhiên$4\,725$.\par\vspace{2pt}
\noindent\textit{Đáp số:} \dotfill\par\vspace{6pt}

\noindent\textbf{Câu 2.} [Thông hiểu] Số La Mã$\mathrm{XVI}$biểu thị số tự nhiên nào trong hệ thập phân?\par\vspace{2pt}
\noindent\textit{Đáp số:} \dotfill\par\vspace{6pt}

\noindent\textbf{Câu 3.} [Thông hiểu] Viết số tự nhiên gồm$5$nghìn,$8$chục và$2$đơn vị dưới dạng một số tự nhiên trong hệ thập phân.\par\vspace{2pt}
\noindent\textit{Đáp số:} \dotfill\par\vspace{6pt}

\noindent\textbf{Câu 4.} [Thông hiểu] Cho số tự nhiên có hai chữ số$\overline{3a}$. Biết chữ số hàng đơn vị là$7$, số đó bằng bao nhiêu?\par\vspace{2pt}
\noindent\textit{Đáp số:} \dotfill\par\vspace{6pt}


% ==============================================================================
% PHẦN IV: TỰ LUẬN (3,0 ĐIỂM)
% ==============================================================================
\vspace{6pt}
\noindent
\textbf{\large PHẦN IV. Tự luận (3,0 điểm)}\\
\textit{\small Thí sinh trình bày chi tiết các bước giải từ câu 1 đến câu 3 vào giấy làm bài.}
\vspace{6pt}
\noindent\textbf{Bài 1 (1,0 điểm).} [Thông hiểu] Cho số tự nhiên$2\,435$:\par\noindent a) Hãy viết số đó thành tổng giá trị các chữ số của nó.\par\noindent b) Cho biết chữ số$4$thuộc hàng nào và có giá trị bằng bao nhiêu?\par\vspace{8pt}

\noindent\textbf{Bài 2 (1,0 điểm).} [Thông hiểu] Thực hiện các yêu cầu sau:\par\noindent a) Viết các số tự nhiên$14$và$26$bằng chữ số La Mã.\par\noindent b) Đọc các số La Mã$\mathrm{XII}$và$\mathrm{XXI}$sang số tự nhiên trong hệ thập phân.\par\vspace{8pt}

\noindent\textbf{Bài 3 (1,0 điểm).} [Thông hiểu] Một số tự nhiên có ba chữ số dạng$\overline{6x3}$, trong đó$x$là chữ số hàng chục.\par\noindent a) Hãy biểu diễn số$\overline{6x3}$thành tổng giá trị các chữ số theo$x$.\par\noindent b) Tìm số đó khi biết$x = 5$và viết lại số đó thành tổng giá trị các chữ số của nó.\par\vspace{8pt}


\begin{center}
\textbf{---------- HẾT ----------}\\
\textit{\small Cán bộ coi thi không giải thích gì thêm.}
\end{center}

\newpage
% ==============================================================================
% HƯỚNG DẪN CHẤM VÀ ĐÁP ÁN CHI TIẾT
% ==============================================================================
\begin{center}
{\Large\bfseries HƯỚNG DẪN CHẤM VÀ ĐÁP ÁN CHI TIẾT}\\[4pt]
{\large\bfseries Toán 6 (KNTT) - Bài 2: Cách ghi số tự nhiên - Bài kiểm tra số 1}\\[2pt]
{\bfseries Mã đề: \texttt{c1\_6\_bai2\_test1}}
\end{center}

\vspace{10pt}
\noindent\textbf{\large 1. ĐÁP ÁN PHẦN I: TRẮC NGHIỆM NHIỀU PHƯƠNG ÁN (3,0 điểm)}
\vspace{4pt}
\begin{center}
\begin{tabular}{|c|c|c|c|c|c|c|c|c|c|c|c|}
\hline
\textbf{1} & \textbf{2} & \textbf{3} & \textbf{4} & \textbf{5} & \textbf{6} & \textbf{7} & \textbf{8} & \textbf{9} & \textbf{10} & \textbf{11} & \textbf{12} \\
\hline
\textbf{B} & \textbf{C} & \textbf{A} & \textbf{C} & \textbf{A} & \textbf{B} & \textbf{B} & \textbf{C} & \textbf{D} & \textbf{A} & \textbf{B} & \textbf{C} \\
\hline
\end{tabular}
\end{center}
\vspace{6pt}


\noindent\textbf{\large 2. ĐÁP ÁN PHẦN II: TRẮC NGHIỆM ĐÚNG / SAI (2,0 điểm)}
\vspace{4pt}
\begin{center}
\begin{tabular}{|c|c|c|c|c|}
\hline
\textbf{Câu hỏi} & \textbf{Ý a)} & \textbf{Ý b)} & \textbf{Ý c)} & \textbf{Ý d)} \\
\hline
\textbf{Câu 1} & \textbf{Đúng} & \textbf{Sai} & \textbf{Đúng} & \textbf{Sai} \\
\hline
\textbf{Câu 2} & \textbf{Đúng} & \textbf{Sai} & \textbf{Đúng} & \textbf{Sai} \\
\hline
\end{tabular}
\end{center}
\vspace{6pt}


\noindent\textbf{\large 3. ĐÁP ÁN PHẦN III: TRẢ LỜI NGẮN (2,0 điểm)}
\vspace{4pt}
\begin{center}
\begin{tabular}{|c|m{11cm}|c|}
\hline
\textbf{Câu} & \textbf{Tóm tắt yêu cầu} & \textbf{Đáp số chính xác} \\
\hline
\textbf{Câu 1} & Tính giá trị của chữ số$7$trong số tự nhiên$4\,725$. & \textbf{700} \\
\hline
\textbf{Câu 2} & Số La Mã$\mathrm{XVI}$biểu thị số tự nhiên nào trong hệ thập phân? & \textbf{16} \\
\hline
\textbf{Câu 3} & Viết số tự nhiên gồm$5$nghìn,$8$chục và$2$đơn vị dưới dạng một số tự nhiên... & \textbf{5082} \\
\hline
\textbf{Câu 4} & Cho số tự nhiên có hai chữ số$\overline{3a}$. Biết chữ số hàng đơn vị là$7$, s... & \textbf{37} \\
\hline
\end{tabular}
\end{center}
\vspace{6pt}


\noindent\textbf{\large 4. LỜI GIẢI CHI TIẾT VÀ BIỂU ĐIỂM TỰ LUẬN (3,0 điểm)}
\vspace{4pt}
\noindent\textbf{\underline{Bài 1 (1,0 điểm):}}\par\vspace{3pt}
Bước 1 (0.5đ): Viết đúng dạng tổng các chữ số:$2\,435 = 2 \times 1\,000 + 4 \times 100 + 3 \times 10 + 5$(hoặc$2\,000 + 400 + 30 + 5$).\par\noindent Bước 2 (0.5đ): Nêu đúng chữ số$4$thuộc hàng trăm (0.25đ) và có giá trị là$400$(0.25đ).\par\vspace{8pt}

\noindent\textbf{\underline{Bài 2 (1,0 điểm):}}\par\vspace{3pt}
Bước 1 (0.5đ): Viết đúng$14 = \mathrm{XIV}$(0.25đ) và$26 = \mathrm{XXVI}$(0.25đ).\par\noindent Bước 2 (0.5đ): Chuyển đổi đúng$\mathrm{XII} = 12$(0.25đ) và$\mathrm{XXI} = 21$(0.25đ).\par\vspace{8pt}

\noindent\textbf{\underline{Bài 3 (1,0 điểm):}}\par\vspace{3pt}
Bước 1 (0.5đ): Biểu diễn đúng$\overline{6x3} = 6 \times 100 + x \times 10 + 3 = 600 + 10 \times x + 3$(hoặc$603 + 10x$).\par\noindent Bước 2 (0.5đ): Thay$x = 5$được số$653$(0.25đ) và biểu diễn thành tổng:$653 = 600 + 50 + 3$(0.25đ).\par\vspace{8pt}

\end{document}
